\documentclass[10pt,a4paper]{amsart}
\usepackage[margin=19mm]{geometry}
\usepackage{amsmath,amssymb,mathtools}
\usepackage[scr=rsfs,cal=euler]{mathalfa}
\usepackage{xfrac}
\usepackage[unicode,hidelinks,bookmarks=false]{hyperref}
\newcommand{\ie}{i.e.\ }
\newcommand{\cf}{cf.\ }
\newcommand{\resp}{respectively\ }
\newcommand{\Subsection}{\subsection}
\providecommand{\nobreakdash}{\nobreak\hbox{-}}
\setlength{\emergencystretch}{2em}
\multlinegap0pt
\usepackage{bm}
\usepackage{bbm}
\usepackage{dsfont}
\usepackage{xspace}
\usepackage{cancel}
\usepackage{mathtools}
\usepackage{amsthm}
\makeatletter
\@ifundefined{theorem}{\newtheorem{theorem}{Theorem}}{}
\@ifundefined{lemma}{\newtheorem{lemma}[theorem]{Lemma}}{}
\makeatother
\DeclareMathOperator{\Z}{\mathbb{Z}}
\newcommand{\gh}{\mathrm{gh}}
\newcommand{\qWmm}{q\text{-}\mathbb{W}_{m'}}
\newcommand{\qWm}{q\text{-}\mathbb{W}_m}
\renewcommand{\ge}{\geqslant}
\renewcommand{\le}{\leqslant}
\title[The cyclosyntomic regulator of a number field]{The cyclosyntomic regulator of a number field}
\author{Tess Bouis, Quentin Gazda}
\date{Source: arXiv:2602.21894v1 (CC BY 4.0)}
\begin{document}
\maketitle

\noindent\emph{Source and licence.} The cyclosyntomic regulator of a number field. Version arXiv:2602.21894v1 (CC BY 4.0), distributed under the Creative Commons Attribution 4.0 International licence, as recorded in the arXiv licence metadata for this exact version and shown on the abstract page https://arxiv.org/abs/2602.21894v1. Full rights notice: licenses/arxiv-2602.21894-CC-BY-4.0.txt.\par\smallskip

\noindent\textit{Editorial scope.} Counted unit: the Lemma whose source label is \texttt{lemmanormpreserveimageofqghostmap} in the article (source0.tex lines 514--517, source label \texttt{lemmanormpreserveimageofqghostmap}), delivered together with the article's own complete original proof of it (source0.tex lines 519--540, one \texttt{proof} environment of 22 lines). No mathematics is written, repaired or re-derived: statement and proof are the article's own text line for line. Required context retained uncounted: lemmaqDwork (lines 482--488). Every definition, display and label of those lines is retained unchanged. Omitted from this package and not redistributed: 1--481, 489--513, 518--518, 541--1338. Nothing inside a retained excerpt is omitted or altered: every retained body is delivered exactly as the article's lines are written, so no omission receipt of any kind accompanies this package. Citations: the retained excerpts carry no citation command at all, so no bibliography entry of the article is transcribed and no reference is redistributed. Reference adaptation: none was needed and none was made; every reference inside the delivered unit resolves inside it, so no unresolved reference or citation is produced. Printed slips kept literal: no printed word, symbol, hypothesis or number of the counted statement or of its proof was altered. Layout and class: the article's own class is \texttt{scrarticle} with packages including inputenc, geometry, fancyhdr, stmaryrd, mathdots, indentfirst, xfrac, amsmath, amscd, amsfonts, amssymb, amsthm, todonotes, eurosym; the wrapper is the lane's pinned \texttt{amsart} editorial wrapper at 10pt on A4, which restates the theorem-like environments the retained text needs and adds the article's own macro definitions that the retained text needs (\texttt{\textbackslash Z}, \texttt{\textbackslash ge}, \texttt{\textbackslash gh}, \texttt{\textbackslash le}, \texttt{\textbackslash qWm}, \texttt{\textbackslash qWmm}), with extra wrapper packages (bm, bbm, dsfont, xspace, cancel, mathtools). The delivered numbering is the unit's own. Assets: no retained span contains a figure environment, a graphics-inclusion command, a PDF-page inclusion, a table or a TikZ picture; no image file is copied into this package, the package declares no asset dependency and no image file is redistributed with it. Licence: version arXiv:2602.21894v1 of this article is distributed under the Creative Commons Attribution 4.0 International licence, as recorded in the arXiv licence metadata for this exact version and shown on the abstract page https://arxiv.org/abs/2602.21894v1. A full rights notice is delivered at licenses/arxiv-2602.21894-CC-BY-4.0.txt. Characters: every retained excerpt is pure ASCII, so the delivered engine, XeLaTeX with its default Latin Modern text font, reports no missing character.
\begin{lemma}[$q$-Dwork's lemma]\label{lemmaqDwork}
    Let $R$ be a commutative ring, and $m \ge 1$ be an integer. Assume that for every prime number $p$ dividing $m$, there exists a ring homomorphism $\varphi_p : R \rightarrow R/p^{v_p(m)}R$ such that $\varphi_p(x) \equiv x^p$ modulo $pR$. Then for every element $(c_e(q))_{e|m} \in \prod_{e|m} R[q]/\Phi_e(q)$, the following are equivalent:
    \begin{enumerate}
        \item $(c_e(q))_{e|m} \in \prod_{e|m} R[q]/\Phi_e(q)$ is in the image of the $q$-ghost map $q$-$\gh$;
        \item there exists a lift $(\tilde{c}_e(q))_{e|m} \in \prod_{e|m} R[q]$ of $(c_e(q))_{e|m} \in \prod_{e|m} R[q]/\Phi_e(q)$ such that, for every prime number $p$ and every integer $e \ge 1$ satisfying $pe|m$, we have $\tilde{c}_e(q) \equiv \varphi_p(\tilde{c}_{pe}(q))$ in $(R/p^{v_p(m/e)})[q]$, where $\varphi_p : R[q] \rightarrow (R/p^{v_p(m/e)})[q]$ is the unique $\Z[q]$-linear ring homomorphism extending $\varphi_p : R \rightarrow R/p^{v_p(m/e)}R$.
    \end{enumerate}
\end{lemma}

\begin{lemma}\label{lemmanormpreserveimageofqghostmap}
    Let $R$ be a commutative ring, and $m,m' \ge 1$ be integers satisfying $m | m'$. Assume that for every prime number $p$ dividing $m$, there exists a ring homomorphism $\varphi_p : R \rightarrow R/p^{v_p(m')}R$ such that $\varphi_p \equiv x^p$ modulo $pR$. Then for every element $c \in \qWm(R)$ with $q$-ghost coordinates $(c_e(q))_{e|m}$,  
    the element $$\Pi_{m'/m}((c_e(q))_{e|m}) := \big(c_{(e,m)}(q^{\frac{e}{(e,m)}})^{\frac{m'}{[e,m]}}\big)_{e|m'} \in \prod_{e|m'} R[q]/\Phi_e(q)$$ is in the image of the $q$-ghost map $q\text{-}\gh : \qWmm(R) \rightarrow \prod_{e | m'} R[q]/\Phi_e(q)$.
\end{lemma}

\begin{proof}
    First note that the morphism of multiplicative monoids $$\Pi_{m'/m} : \prod_{e|m} R[q]/\Phi_e(q) \rightarrow \prod_{e|m'} R[q]/\Phi_e(q)$$ is well-defined, because $\Phi_e(q)$ divides $\Phi_{(e,m)}(q^{\frac{e}{(e,m)}})$ in the commutative ring $\Z[q]$. Moreover, for any integers $m,m',m'' \ge 1$ satisfying $m|m'|m''$, we have the equality $$\Pi_{m''/m'} \circ \Pi_{m'/m} = \Pi_{m''/m}.$$ Indeed, this is a consequence of the series of equalities
    $$c_{((e,m'),m)}\big(q^{\frac{(e,m')}{((e,m'),m)}\cdot\frac{e}{(e,m')}}\big)^{\frac{m'}{[(e,m'),m]}\cdot\frac{m''}{[e,m']}} = c_{(e,m)}(q^{\frac{e}{(e,m)}})^{\frac{m'}{[(e,m'),m]}\cdot\frac{m''}{[e,m']}} = c_{(e,m)}(q^{\frac{e}{(e,m)}})^{\frac{m''}{[e,m]}}$$
    where the equality $\tfrac{m'}{[(e,m'),m]\cdot[e,m']}=\tfrac{1}{[e,m]}$ follows from the fact that for any integers $a,b,c \ge 0$ satisfying $c \le b$, we have the equality 
    $b-\max(\min(a,b),c)-\max(a,b)=-\max(a,c)$. 
    
    In particular, it suffices to prove the desired claim when $m'/m$ is a prime number, which we assume from now on.
    
    Let $(c_e(q))_{e|m} \in \prod_{e|m} R[q]/\Phi_e(q)$ be an element in the image of the $q$-ghost map. We use the $q$-Dwork lemma (Lemma~\ref{lemmaqDwork}) to prove that $\Pi_{m'/m}((c_e(q))_{e|m}) \in \prod_{e|m'} R[q]/\Phi_e(q)$ is in the image of the $q$-ghost map. Let $(\tilde{c}_e(q))_{e|m} \in \prod_{e|m} R[q]$ be a lift of $(c_e(q))_{e|m} \in \prod_{e|m} R[q]/\Phi_e(q)$ such that, for every prime number $p$ and every integer $e \ge 1$ satisfying $pe|m$, we have $\tilde{c}_e(q) \equiv \varphi_p(\tilde{c}_{pe}(q))$ in $(R/p^{v_p(m/e)})[q]$. It then suffices to prove that, for every prime number $p$ and every integer $e \ge 1$ satisfying $pe|m'$, we have $\varphi_p\big(\tilde{c}_{(pe,m)}(q^{\frac{pe}{(pe,m)}})^{\frac{m'}{[pe,m]}}\big) \equiv \tilde{c}_{(e,m)}(q^{\frac{e}{(e,m)}})^{\frac{m'}{[e,m]}}$ in $(R/p^{v_p(m'/e)})[q]$. We distinguish several cases.
    
    Assume first that $v_p(e)>v_p(m)$. Because $m'/m$ is a prime number, this implies that $m'/m=p$, so the condition $pe|m'$ can not be satisfied in this case,  
    and the desired statement is vacuously true.

    Assume now that $v_p(e)=v_p(m)$. If $v_p(e)=v_p(m')$, then $v_p(m'/e)=0$ and the desired statement is again vacuously true, so we assume that $v_p(e)=v_p(m)<v_p(m')$. Because $m'/m$ is a prime number, this implies that $m'=pm$ and that $v_p(m'/e)=1$. The desired statement then follows from the series of equalities
    $$\varphi_p\big(\tilde{c}_{(pe,m)}(q^{\frac{pe}{(pe,m)}})^{\frac{m'}{[pe,m]}}\big) = \varphi_p\big(\tilde{c}_{(e,m)}(q^{\frac{pe}{(e,m)}})\big)^{\frac{m'}{p\cdot[e,m]}} \equiv \tilde{c}_{(e,m)}(q^{\frac{e}{(e,m)}})^{p\cdot\frac{m'}{p\cdot [e,m]}} = \tilde{c}_{(e,m)}(q^{\frac{e}{(e,m)}})^{\frac{m'}{[e,m]}}$$
    where the congruence holds modulo $p$, using that $\varphi_p : R \rightarrow R/p^{v_p(m')}R$ is a lift of Frobenius. 

    Assume finally that $v_p(e)<v_p(m)$. In this case, the desired statement follows similarly from the series of equalities
    $$\varphi_p\big(\tilde{c}_{(pe,m)}(q^{\frac{pe}{(pe,m)}})^{\frac{m'}{[pe,m]}}\big) = \varphi_p\big(\tilde{c}_{p\cdot(e,m)}(q^{\frac{pe}{p\cdot(e,m)}})\big)^{\frac{m'}{[e,m]}} \equiv \tilde{c}_{(e,m)}(q^{\frac{e}{(e,m)}})^{\frac{m'}{[e,m]}} = \tilde{c}_{(e,m)}(q^{\frac{e}{(e,m)}})^{\frac{m'}{[e,m]}}$$
    where the congruence holds modulo $p^{v_p(m'/e)}$, as a consequence of the assumption on $(\tilde{c}_e(q))_{e|m}$.  
    More precisely, either $m'/m$ is different from $p$, in which case $v_p(m'/e)=v_p(m/(e,m))$ and this is indeed a consequence of the assumption, or $m'/m=p$, in which case we use the fact that $a \equiv b$ modulo $p^{v_p(m/e)}$ implies $a^p \equiv b^p$ modulo $p^{v_p(m'/e)}$. 
\end{proof}

\end{document}
